\documentclass[a4paper,10pt]{article}
\usepackage[margin=1.5cm]{geometry}
\usepackage{multicol}
\usepackage{hyperref}
\usepackage{enumitem}
\usepackage{listings}
\usepackage{graphicx}


\setlength{\parindent}{0pt}

\begin{document}


\begin{center}
    {\Large \texttt{Modelo TCP/IP}}\\
\end{center}

\tableofcontents
\newpage



\section{Modelo TCP/IP}
Al igual que el modelo OSI, el \textbf{modelo TCP/IP} propone el conjunto de protocolos más usado de comunicación que permite a los programas de dispositivos informáticos intercambiar
mensajes a través de la internet. Está diseñado para \textbf{enviar paquetes a través de la red y garantizar la entrega exitosa de datos y mensajes entre dispositivos}.


\section{Capas del modelo TCP/IP}
El modelo TCP/IP se conforma de 4 capas de protocolos:

\begin{itemize}
    \item Capa 4 - \textbf{Aplicación},
    \item Capa 3 - \textbf{Transporte},
    \item Capa 2 - \textbf{Red} e
    \item Capa 1 - \textbf{Interfáz de Red}.
\end{itemize}


\subsection{Capa 4 - Aplicación}
La \textbf{capa de Aplicación} incluye los protocolos utilizados por la mayoría de las aplicaciones para \textbf{proporcionar servicios de usuario o intercambiar datos de aplicaciones} 
a través de las conexiones de red establecidas por los protocolos de las capas inferiores. Esto puede incluir algunos servicios básicos de soporte de red, como protocolos de enrutamiento y configuración de host. 

\

Algunos protocolos presentes en esta capa son:

\begin{itemize}
    \item Hypertext Transfer Protocol (Secure) - HTTP(S).
    \item File Transfer Protocol - FTP.
    \item Dinamic Host Config Protocol - DHCP.
    \item Simple Mail Transfer Protoco - SMTP.
    \item Domain Name Service - DNS.
\end{itemize}

\begin{center}
    \textit{Aquí van los protocolos que solicitan y reciben los recursos en su formato final.}
\end{center}


Respecto al \textbf{modelo OSI}, la capa de aplicación en el modelo TCP/IP corresponde a una combinación de la quinta (\textbf{sesión}), sexta (\textbf{presentación}) y séptima capa (\textbf{aplicación}).



\subsection{Capa 3 - Transporte}
En la \textbf{Capa de Transporte} se 


\subsection{Capa 2 - Red}

\subsection{Capa 1 - Interfáz de Red}



\subsection{Modelo OSI vs TCP/IP}

\section{TCP}


\section{IP}



\end{document}